\chapter{Metodología} \label{sec:Metodos}

\begin{itemize}
    \item Tipo de investigación
    \item Diseño experimental o metodológico
    \item Fuentes y técnicas de recolección de datos
    \item Procesamiento y análisis
    \item Consideraciones éticas

\end{itemize}


En este capítulo se presenta la metodología de trabajo propuesta para darle solución al problema de investigación. Describe los modelos o algoritmos desarrollados, las técnicas y  métodos empleados, los materiales, software e insumos usados en la investigación.

En Overleaf, una plataforma en línea para escribir documentos LaTeX colaborativamente, es sencillo incluir citas y referencias utilizando BibTeX. Aquí te explico cómo hacerlo. En primer lugar, debes crear un archivo con extensión .bib en tu proyecto de Overleaf. Este archivo contendrá todas las referencias bibliográficas que desees citar en tu documento LaTeX. 

Luego, en ese archivo .bib, debes agregar las entradas de las referencias bibliográficas que quieras citar. Cada entrada tiene una estructura específica que varía según el tipo de referencia, como libro, artículo, tesis, entre otros. Se pueden obtener estas entradas desde muchas fuentes, Mendeley, Zotero, o incluso ScholarGoogle. Buscas tus referencias y las pegas en el archivo .bib. 

Luego citar en muy sencillo, basta con hacer algo como \cite{liakos2018machine}.